\section{Discussion}

\emph{What the composite hides.} The composite can hide two different problems, and each needs its own audit. When a dimension carries signal, the composite detects it worse than the model's own output for that dimension (positive $\Delta_d$); showing the per-dimension outputs next to the composite removes this loss at no cost. When federated training fails to learn a dimension at all (Sleep and Mental in the baseline), $\Delta_d$ is small because there is little left to hide, and only a per-dimension AUC audit catches the failure. Comparing per-client MG with a fixed cut-point catches neither.

\emph{Clinical relevance.} All labels except Diet use published cut-points (Table~\ref{tab:dims}). An AUC of 0.5 means no discrimination, and under one commonly cited descriptive convention 0.7--0.8 is termed acceptable~\cite{mandrekar2010}; these bands are not clinical-utility thresholds. Validated screeners such as the PHQ-9 do far better~\cite{kroenke2001}. Most per-dimension AUCs here are below 0.7 and none exceeds 0.8 when averaged over seeds, so none of these models is fit for screening individuals, masked or not. We therefore present MG as an \emph{audit} statistic whose severity is judged against sampling noise, not as a clinical-utility threshold; no guideline defines one for composite health scores.

\emph{Recommendations for auditing composite health models under FL.}
(i) Report MG with its chance level (MG-excess and null percentile), never raw. (ii) Plan per-client test sets with Eq.~(\ref{eq:closed}) (here about 300 records to resolve a gap of 0.10); audit smaller clients pooled or grouped. (iii) Treat AUCs with fewer than 5 cases of either class as missing, not as 0.5. (iv) Check each dimension's population AUC against chance and against a centralized model; an output at chance is a failure whatever the MG. (v) Report sensitivity at a validation-tuned threshold; at 0.5 it mostly reflects prevalence. (vi) Report $\Delta_d$ to show how much the composite hides.

\emph{Limitations.} NHANES is cross-sectional survey and exam data, not continuous wearable streams, so AUCs are modest and the results describe a setting with little information per person. The federations are simulated (three partitions of one survey, synchronous rounds), and survey weights are not applied. The random-score null treats dimensions as independent; the label-permutation null, used for the full protocol, keeps their dependence. Ditto's selected $\lambda$ was the largest value tried (1). Diet uses a sample-median split.
